% =================================================
% 2. ENVIRONMENT ANALYSIS
% Included from main.tex with % =================================================
% 2. ENVIRONMENT ANALYSIS
% Included from main.tex with % =================================================
% 2. ENVIRONMENT ANALYSIS
% Included from main.tex with % =================================================
% 2. ENVIRONMENT ANALYSIS
% Included from main.tex with \input{section2_environment}
% =================================================
\section{Environment Analysis}
\label{sec:environment}

Before choosing the agents, we looked at the environment they will work in, using the five
properties of Lecture~2 \cite{lecture2, wooldridge2009}. It is important to keep in mind
that our agents do not see a real airspace. They see one scenario, given as files: the
three required by the assignment (airspace graph, aircraft and flight plans) and the three
optional ones we added (weather, restricted zones and airport constraints, see
Table~\ref{tab:inputs}). The weather comes from Open-Meteo \cite{openmeteo_docs}, but we
download it once with a script and save it in the scenario folder, so the system never
downloads anything while it runs. The system reads the files, writes the report and stops
\cite{mas_assignment}. Most of our answers below come from this, and for a real ATC centre
almost all of them would be the opposite.

\subsection{Accessible vs. Inaccessible}

We consider the environment \textbf{accessible}. All the information the agents need is in
the input files from the beginning: waypoints, distances, speeds, routes, departure times,
cruise levels, weather and restrictions. Data Management normalises it once, and after that
every crew reads the same intermediate files, so no agent has a partial or old version of
the scenario.

The only limitation is the weather. The snapshot is a forecast for one moment, so it is not
exactly what the atmosphere will be like. But for our system it is the complete
information, and there is nothing else to look at. A real controller has it much harder:
radar data has noise and delays, and the controller does not always know what the pilot is
going to do.

\subsection{Deterministic vs. Non-deterministic}

Inside our model the environment is \textbf{deterministic}. Aircraft follow their flight
plans exactly, the wind in the snapshot does not change, and the constraint zones are
created from fixed thresholds. If we run the same scenario and apply the same manoeuvre
(level change, speed change, delay or hold, or reroute), we always get the same
trajectories and the same conflicts. For this reason all the calculations are done by
deterministic tools and not by the language model (Section~\ref{sec:extra-crews}).

The LLM agents can give different answers in different runs, but we see this as a property
of the agents, not of the environment. In the real world it is different: the wind changes,
storms move, and a pilot may take some time to follow an instruction, so the same
instruction does not always have the same effect.

\subsection{Episodic vs. Non-episodic}

If we compare runs, the environment is \textbf{episodic}. Each scenario is solved on its own
and nothing is kept for the next one. Inside one run it is not so simple, because the
decisions depend on each other: if we move one aircraft to solve a conflict, it may end up
in a new conflict with a third aircraft, or inside a constraint zone. This is exactly why we
added the Safety Validation crew. It checks the new flight plans again and, if a manoeuvre
is rejected, sends it back to Tactical Resolution (at most twice).

\subsection{Static vs. Dynamic}

The environment is \textbf{static}. During a run no new flights, positions, weather or
restrictions arrive, and the scenario only changes because of the manoeuvres our own agents
propose. Restricted zones and airport constraints have time windows when they are active,
but we know these windows from the start, so they are part of the input and not a change in
the environment.

The escalation to a human controller does not change this either. When a conflict is still
rejected after two iterations, it is marked in the report as ``requires human controller'',
but the system does not stop to wait for an answer. A real airspace is very dynamic, and a
real ATC system would have to replan all the time.

\subsection{Discrete vs. Continuous}

The real problem is \textbf{continuous}: positions, altitudes and time can take any value.
However, our system works with it in a \textbf{discrete} way. The airspace is a graph with a
finite number of waypoints, the inputs are a finite list of records, the detector samples
the trajectories at fixed time steps, each constraint zone has an area, an altitude band and
a time window, and Tactical Resolution can only choose between four types of manoeuvre. The
values of a manoeuvre (for example, the new flight level) are still numbers, so we would say
it is a discrete model of a continuous world.

\subsection{Other Properties}

Russell and Norvig also describe some properties that are not in Lecture~2
\cite{russellnorvig2020}. Two of them help to describe our case. The environment is
\textbf{multi-agent}, but all the agents cooperate, since they share the same goal (a safe
and complete report). The aircraft are not agents, because they only follow their flight
plans. The environment is also \textbf{known}, because the rules, such as the separation
minima and the thresholds of the zones, are fixed before the run starts.

\subsection{Environment Summary Table}

Table~\ref{tab:environment} summarises our classification and how each property affected
the design.

\begin{table}[H]
\centering
\small
\caption{Classification of the environment and its effect on the design.}
\label{tab:environment}
\begin{tabularx}{\textwidth}{@{}l l X@{}}
\toprule
\textbf{Property} & \textbf{Classification} & \textbf{Reason and effect on the design} \\
\midrule
Accessibility & Accessible & All the data is in the input files, including the stored
  weather. All crews read the same intermediate files. \\
Determinism & Deterministic (in the model) & Same scenario and same manoeuvre give the same
  result, so calculations are done by tools. \\
Episodicity & Episodic between runs & Decisions inside a run affect each other, so we added
  the Safety Validation loop. \\
Dynamism & Static & Nothing new arrives during the run, so the system runs once and does
  not need to replan. \\
Continuity & Discrete (by design) & Waypoint graph, sampled trajectories, constraint zones
  and four types of manoeuvre. \\
\bottomrule
\end{tabularx}
\end{table}

% =================================================
\section{Environment Analysis}
\label{sec:environment}

Before choosing the agents, we looked at the environment they will work in, using the five
properties of Lecture~2 \cite{lecture2, wooldridge2009}. It is important to keep in mind
that our agents do not see a real airspace. They see one scenario, given as files: the
three required by the assignment (airspace graph, aircraft and flight plans) and the three
optional ones we added (weather, restricted zones and airport constraints, see
Table~\ref{tab:inputs}). The weather comes from Open-Meteo \cite{openmeteo_docs}, but we
download it once with a script and save it in the scenario folder, so the system never
downloads anything while it runs. The system reads the files, writes the report and stops
\cite{mas_assignment}. Most of our answers below come from this, and for a real ATC centre
almost all of them would be the opposite.

\subsection{Accessible vs. Inaccessible}

We consider the environment \textbf{accessible}. All the information the agents need is in
the input files from the beginning: waypoints, distances, speeds, routes, departure times,
cruise levels, weather and restrictions. Data Management normalises it once, and after that
every crew reads the same intermediate files, so no agent has a partial or old version of
the scenario.

The only limitation is the weather. The snapshot is a forecast for one moment, so it is not
exactly what the atmosphere will be like. But for our system it is the complete
information, and there is nothing else to look at. A real controller has it much harder:
radar data has noise and delays, and the controller does not always know what the pilot is
going to do.

\subsection{Deterministic vs. Non-deterministic}

Inside our model the environment is \textbf{deterministic}. Aircraft follow their flight
plans exactly, the wind in the snapshot does not change, and the constraint zones are
created from fixed thresholds. If we run the same scenario and apply the same manoeuvre
(level change, speed change, delay or hold, or reroute), we always get the same
trajectories and the same conflicts. For this reason all the calculations are done by
deterministic tools and not by the language model (Section~\ref{sec:extra-crews}).

The LLM agents can give different answers in different runs, but we see this as a property
of the agents, not of the environment. In the real world it is different: the wind changes,
storms move, and a pilot may take some time to follow an instruction, so the same
instruction does not always have the same effect.

\subsection{Episodic vs. Non-episodic}

If we compare runs, the environment is \textbf{episodic}. Each scenario is solved on its own
and nothing is kept for the next one. Inside one run it is not so simple, because the
decisions depend on each other: if we move one aircraft to solve a conflict, it may end up
in a new conflict with a third aircraft, or inside a constraint zone. This is exactly why we
added the Safety Validation crew. It checks the new flight plans again and, if a manoeuvre
is rejected, sends it back to Tactical Resolution (at most twice).

\subsection{Static vs. Dynamic}

The environment is \textbf{static}. During a run no new flights, positions, weather or
restrictions arrive, and the scenario only changes because of the manoeuvres our own agents
propose. Restricted zones and airport constraints have time windows when they are active,
but we know these windows from the start, so they are part of the input and not a change in
the environment.

The escalation to a human controller does not change this either. When a conflict is still
rejected after two iterations, it is marked in the report as ``requires human controller'',
but the system does not stop to wait for an answer. A real airspace is very dynamic, and a
real ATC system would have to replan all the time.

\subsection{Discrete vs. Continuous}

The real problem is \textbf{continuous}: positions, altitudes and time can take any value.
However, our system works with it in a \textbf{discrete} way. The airspace is a graph with a
finite number of waypoints, the inputs are a finite list of records, the detector samples
the trajectories at fixed time steps, each constraint zone has an area, an altitude band and
a time window, and Tactical Resolution can only choose between four types of manoeuvre. The
values of a manoeuvre (for example, the new flight level) are still numbers, so we would say
it is a discrete model of a continuous world.

\subsection{Other Properties}

Russell and Norvig also describe some properties that are not in Lecture~2
\cite{russellnorvig2020}. Two of them help to describe our case. The environment is
\textbf{multi-agent}, but all the agents cooperate, since they share the same goal (a safe
and complete report). The aircraft are not agents, because they only follow their flight
plans. The environment is also \textbf{known}, because the rules, such as the separation
minima and the thresholds of the zones, are fixed before the run starts.

\subsection{Environment Summary Table}

Table~\ref{tab:environment} summarises our classification and how each property affected
the design.

\begin{table}[H]
\centering
\small
\caption{Classification of the environment and its effect on the design.}
\label{tab:environment}
\begin{tabularx}{\textwidth}{@{}l l X@{}}
\toprule
\textbf{Property} & \textbf{Classification} & \textbf{Reason and effect on the design} \\
\midrule
Accessibility & Accessible & All the data is in the input files, including the stored
  weather. All crews read the same intermediate files. \\
Determinism & Deterministic (in the model) & Same scenario and same manoeuvre give the same
  result, so calculations are done by tools. \\
Episodicity & Episodic between runs & Decisions inside a run affect each other, so we added
  the Safety Validation loop. \\
Dynamism & Static & Nothing new arrives during the run, so the system runs once and does
  not need to replan. \\
Continuity & Discrete (by design) & Waypoint graph, sampled trajectories, constraint zones
  and four types of manoeuvre. \\
\bottomrule
\end{tabularx}
\end{table}

% =================================================
\section{Environment Analysis}
\label{sec:environment}

Before choosing the agents, we looked at the environment they will work in, using the five
properties of Lecture~2 \cite{lecture2, wooldridge2009}. It is important to keep in mind
that our agents do not see a real airspace. They see one scenario, given as files: the
three required by the assignment (airspace graph, aircraft and flight plans) and the three
optional ones we added (weather, restricted zones and airport constraints, see
Table~\ref{tab:inputs}). The weather comes from Open-Meteo \cite{openmeteo_docs}, but we
download it once with a script and save it in the scenario folder, so the system never
downloads anything while it runs. The system reads the files, writes the report and stops
\cite{mas_assignment}. Most of our answers below come from this, and for a real ATC centre
almost all of them would be the opposite.

\subsection{Accessible vs. Inaccessible}

We consider the environment \textbf{accessible}. All the information the agents need is in
the input files from the beginning: waypoints, distances, speeds, routes, departure times,
cruise levels, weather and restrictions. Data Management normalises it once, and after that
every crew reads the same intermediate files, so no agent has a partial or old version of
the scenario.

The only limitation is the weather. The snapshot is a forecast for one moment, so it is not
exactly what the atmosphere will be like. But for our system it is the complete
information, and there is nothing else to look at. A real controller has it much harder:
radar data has noise and delays, and the controller does not always know what the pilot is
going to do.

\subsection{Deterministic vs. Non-deterministic}

Inside our model the environment is \textbf{deterministic}. Aircraft follow their flight
plans exactly, the wind in the snapshot does not change, and the constraint zones are
created from fixed thresholds. If we run the same scenario and apply the same manoeuvre
(level change, speed change, delay or hold, or reroute), we always get the same
trajectories and the same conflicts. For this reason all the calculations are done by
deterministic tools and not by the language model (Section~\ref{sec:extra-crews}).

The LLM agents can give different answers in different runs, but we see this as a property
of the agents, not of the environment. In the real world it is different: the wind changes,
storms move, and a pilot may take some time to follow an instruction, so the same
instruction does not always have the same effect.

\subsection{Episodic vs. Non-episodic}

If we compare runs, the environment is \textbf{episodic}. Each scenario is solved on its own
and nothing is kept for the next one. Inside one run it is not so simple, because the
decisions depend on each other: if we move one aircraft to solve a conflict, it may end up
in a new conflict with a third aircraft, or inside a constraint zone. This is exactly why we
added the Safety Validation crew. It checks the new flight plans again and, if a manoeuvre
is rejected, sends it back to Tactical Resolution (at most twice).

\subsection{Static vs. Dynamic}

The environment is \textbf{static}. During a run no new flights, positions, weather or
restrictions arrive, and the scenario only changes because of the manoeuvres our own agents
propose. Restricted zones and airport constraints have time windows when they are active,
but we know these windows from the start, so they are part of the input and not a change in
the environment.

The escalation to a human controller does not change this either. When a conflict is still
rejected after two iterations, it is marked in the report as ``requires human controller'',
but the system does not stop to wait for an answer. A real airspace is very dynamic, and a
real ATC system would have to replan all the time.

\subsection{Discrete vs. Continuous}

The real problem is \textbf{continuous}: positions, altitudes and time can take any value.
However, our system works with it in a \textbf{discrete} way. The airspace is a graph with a
finite number of waypoints, the inputs are a finite list of records, the detector samples
the trajectories at fixed time steps, each constraint zone has an area, an altitude band and
a time window, and Tactical Resolution can only choose between four types of manoeuvre. The
values of a manoeuvre (for example, the new flight level) are still numbers, so we would say
it is a discrete model of a continuous world.

\subsection{Other Properties}

Russell and Norvig also describe some properties that are not in Lecture~2
\cite{russellnorvig2020}. Two of them help to describe our case. The environment is
\textbf{multi-agent}, but all the agents cooperate, since they share the same goal (a safe
and complete report). The aircraft are not agents, because they only follow their flight
plans. The environment is also \textbf{known}, because the rules, such as the separation
minima and the thresholds of the zones, are fixed before the run starts.

\subsection{Environment Summary Table}

Table~\ref{tab:environment} summarises our classification and how each property affected
the design.

\begin{table}[H]
\centering
\small
\caption{Classification of the environment and its effect on the design.}
\label{tab:environment}
\begin{tabularx}{\textwidth}{@{}l l X@{}}
\toprule
\textbf{Property} & \textbf{Classification} & \textbf{Reason and effect on the design} \\
\midrule
Accessibility & Accessible & All the data is in the input files, including the stored
  weather. All crews read the same intermediate files. \\
Determinism & Deterministic (in the model) & Same scenario and same manoeuvre give the same
  result, so calculations are done by tools. \\
Episodicity & Episodic between runs & Decisions inside a run affect each other, so we added
  the Safety Validation loop. \\
Dynamism & Static & Nothing new arrives during the run, so the system runs once and does
  not need to replan. \\
Continuity & Discrete (by design) & Waypoint graph, sampled trajectories, constraint zones
  and four types of manoeuvre. \\
\bottomrule
\end{tabularx}
\end{table}

% =================================================
\section{Environment Analysis}
\label{sec:environment}

Before choosing the agents, we looked at the environment they will work in, using the five
properties of Lecture~2 \cite{lecture2, wooldridge2009}. It is important to keep in mind
that our agents do not see a real airspace. They see one scenario, given as files: the
three required by the assignment (airspace graph, aircraft and flight plans) and the three
optional ones we added (weather, restricted zones and airport constraints, see
Table~\ref{tab:inputs}). The weather comes from Open-Meteo \cite{openmeteo_docs}, but we
download it once with a script and save it in the scenario folder, so the system never
downloads anything while it runs. The system reads the files, writes the report and stops
\cite{mas_assignment}. Most of our answers below come from this, and for a real ATC centre
almost all of them would be the opposite.

\subsection{Accessible vs. Inaccessible}

We consider the environment \textbf{accessible}. All the information the agents need is in
the input files from the beginning: waypoints, distances, speeds, routes, departure times,
cruise levels, weather and restrictions. Data Management normalises it once, and after that
every crew reads the same intermediate files, so no agent has a partial or old version of
the scenario.

The only limitation is the weather. The snapshot is a forecast for one moment, so it is not
exactly what the atmosphere will be like. But for our system it is the complete
information, and there is nothing else to look at. A real controller has it much harder:
radar data has noise and delays, and the controller does not always know what the pilot is
going to do.

\subsection{Deterministic vs. Non-deterministic}

Inside our model the environment is \textbf{deterministic}. Aircraft follow their flight
plans exactly, the wind in the snapshot does not change, and the constraint zones are
created from fixed thresholds. If we run the same scenario and apply the same manoeuvre
(level change, speed change, delay or hold, or reroute), we always get the same
trajectories and the same conflicts. For this reason all the calculations are done by
deterministic tools and not by the language model (Section~\ref{sec:extra-crews}).

The LLM agents can give different answers in different runs, but we see this as a property
of the agents, not of the environment. In the real world it is different: the wind changes,
storms move, and a pilot may take some time to follow an instruction, so the same
instruction does not always have the same effect.

\subsection{Episodic vs. Non-episodic}

If we compare runs, the environment is \textbf{episodic}. Each scenario is solved on its own
and nothing is kept for the next one. Inside one run it is not so simple, because the
decisions depend on each other: if we move one aircraft to solve a conflict, it may end up
in a new conflict with a third aircraft, or inside a constraint zone. This is exactly why we
added the Safety Validation crew. It checks the new flight plans again and, if a manoeuvre
is rejected, sends it back to Tactical Resolution (at most twice).

\subsection{Static vs. Dynamic}

The environment is \textbf{static}. During a run no new flights, positions, weather or
restrictions arrive, and the scenario only changes because of the manoeuvres our own agents
propose. Restricted zones and airport constraints have time windows when they are active,
but we know these windows from the start, so they are part of the input and not a change in
the environment.

The escalation to a human controller does not change this either. When a conflict is still
rejected after two iterations, it is marked in the report as ``requires human controller'',
but the system does not stop to wait for an answer. A real airspace is very dynamic, and a
real ATC system would have to replan all the time.

\subsection{Discrete vs. Continuous}

The real problem is \textbf{continuous}: positions, altitudes and time can take any value.
However, our system works with it in a \textbf{discrete} way. The airspace is a graph with a
finite number of waypoints, the inputs are a finite list of records, the detector samples
the trajectories at fixed time steps, each constraint zone has an area, an altitude band and
a time window, and Tactical Resolution can only choose between four types of manoeuvre. The
values of a manoeuvre (for example, the new flight level) are still numbers, so we would say
it is a discrete model of a continuous world.

\subsection{Other Properties}

Russell and Norvig also describe some properties that are not in Lecture~2
\cite{russellnorvig2020}. Two of them help to describe our case. The environment is
\textbf{multi-agent}, but all the agents cooperate, since they share the same goal (a safe
and complete report). The aircraft are not agents, because they only follow their flight
plans. The environment is also \textbf{known}, because the rules, such as the separation
minima and the thresholds of the zones, are fixed before the run starts.

\subsection{Environment Summary Table}

Table~\ref{tab:environment} summarises our classification and how each property affected
the design.

\begin{table}[H]
\centering
\small
\caption{Classification of the environment and its effect on the design.}
\label{tab:environment}
\begin{tabularx}{\textwidth}{@{}l l X@{}}
\toprule
\textbf{Property} & \textbf{Classification} & \textbf{Reason and effect on the design} \\
\midrule
Accessibility & Accessible & All the data is in the input files, including the stored
  weather. All crews read the same intermediate files. \\
Determinism & Deterministic (in the model) & Same scenario and same manoeuvre give the same
  result, so calculations are done by tools. \\
Episodicity & Episodic between runs & Decisions inside a run affect each other, so we added
  the Safety Validation loop. \\
Dynamism & Static & Nothing new arrives during the run, so the system runs once and does
  not need to replan. \\
Continuity & Discrete (by design) & Waypoint graph, sampled trajectories, constraint zones
  and four types of manoeuvre. \\
\bottomrule
\end{tabularx}
\end{table}
